\chapter{联邦学习}
\begin{quote}\small
\textit{本章摘要。}本章把“数据不能集中”作为约束，重新审视训练、通信、隐私和公平性。先定义联邦优化与非独立同分布，再分析 Federated Averaging (FedAvg)、本地更新、差分隐私、安全聚合和个性化方法，最后回到系统延迟、参与方差异和攻击面。联邦学习解决的是数据协作边界，不会自动解决统计偏差。
\end{quote}
上一章主动学习讨论把有限标注预算花在哪些样本上，隐含的问题是所选样本能否进入同一个学习流程。本章进一步撤去数据可集中这一便利条件：不同工厂或设备群既可能具有不同工况，也可能不能把原始记录汇集到同一处。故障预测目标没有改变，改变的是谁持有数据、谁计算梯度，以及谁能够观察更新。

假设几家工厂不能交换原始记录，但希望一起改进故障模型。首先要决定大工厂是否因样本多而占更大权重，再看各厂独立训练多步以后，更新为什么可能朝不同方向走。通信中断、设备掉线和更新中泄露的信息，也都要考虑。本章从加权经验风险推导 FedAvg，再解释近端项、控制变量和个性化方法。评价时既看平均成绩，也看表现最差的工厂，并记录通信成本和隐私预算。原始数据留在本地，并不自动保证上传的更新安全。
\section{问题定义}

% auxiliary-resource-guide
本章末的“辅助学习材料”列出与核心方法对应的讲解和实践入口，可在阅读相关推导后，按所列小节对照学习。

\begin{figure}[htbp]
\centering
\begin{tikzpicture}[node distance=0.8cm and 1.0cm]
 \node[idi state, minimum width=2.4cm] (server) {服务器\\全局参数 $\theta^t$};
 \node[idi input, below left=1.0cm and 1.0cm of server, minimum width=2.0cm] (c1) {工厂 1\\本地数据};
 \node[idi input, below=1.0cm of server, minimum width=2.0cm] (c2) {工厂 2\\本地数据};
 \node[idi input, below right=1.0cm and 1.0cm of server, minimum width=2.0cm] (c3) {工厂 $K$\\本地数据};
 \draw[idi arrow,<->] (server)--(c1); \draw[idi arrow,<->] (server)--(c2); \draw[idi arrow,<->] (server)--(c3);
 \node[below=0.3cm of c2,font=\scriptsize] {双向通信：下发模型、回传更新；服务器按权重聚合};
\end{tikzpicture}
\caption{联邦学习的基本循环。原始数据留在客户端，服务器只接收模型参数或更新。}
\label{fig:federated-basic-loop}
\end{figure}
图\ref{fig:federated-basic-loop}区分本地数据路径与服务器聚合路径；客户端在原始记录上训练，服务器通过汇总更新协调模型，而不是集中收集训练样本。
FedAvg 通过客户端本地更新与服务器加权聚合实现分散训练\cite{mcmahan2017fedavg}。Federated Proximal Optimization (FedProx) 加入近端约束，Stochastic Controlled Averaging for Federated Learning (SCAFFOLD) 使用控制变量校正客户端漂移\cite{li2020fedprox,karimireddy2020scaffold}；它们改变优化机制，并不单独提供隐私保证。
联邦学习在多个客户端保留本地数据，仅交换模型更新。客户端$k$有目标$F_k(\theta)$和权重$p_k$，全局目标为
\begin{equation}
 \min_\theta F(\theta)=\sum_{k=1}^{K}p_kF_k(\theta).
\end{equation}
FedAvg在每轮下发全局参数，客户端执行若干步本地 SGD，再按样本量聚合：
\begin{equation}
 \theta^{t+1}=\sum_{k=1}^{K}p_k\theta_k^{t+1}.
\end{equation}

\subsection{权重、抽样与集中式基线}
设参数$\theta\in\R^d$，客户端$k$持有$n_k$个样本，且
\begin{equation}
 F_k(\theta)=\frac1{n_k}\sum_{i=1}^{n_k}\ell(\theta;x_{ki},y_{ki}),
 \qquad p_k\geq0,\quad\sum_kp_k=1.
\end{equation}
取$p_k=n_k/\sum_jn_j$得到所有样本等权的经验风险；取$p_k=1/K$得到所有工厂等权的风险。二者不是实现上的可互换选项：一家拥有大量正常记录的大工厂，在第一种目标中可能压过数据稀少但故障代价高的小工厂。权重应先由业务目标定义，而不是由方便上传的数据量事后决定。

若所有客户端参与且只做一步全梯度下降，则
\begin{align}
 \sum_kp_k(\theta^t-\eta\nabla F_k(\theta^t))
 &=\theta^t-\eta\sum_kp_k\nabla F_k(\theta^t)\\
 &=\theta^t-\eta\nabla F(\theta^t).\label{eq:federated-centralized-equivalence}
\end{align}
式\eqref{eq:federated-centralized-equivalence}表明，这个特例与集中式全梯度更新严格相同。若使用随机梯度$g_k$，需要条件无偏性$\E[g_k\mid\theta^t]=\nabla F_k(\theta^t)$才有期望意义上的相同；模型中的随机状态、归一化统计量等还应采用相容的处理方式。

部分参与使权重问题更复杂。令$I_k^t\in\{0,1\}$表示是否参与，已知参与概率$\pi_k>0$，则在更新与抽样条件相容时，
\begin{equation}
 \widehat g^t=\sum_k I_k^t\frac{p_k}{\pi_k}g_k^t,
 \qquad \E[\widehat g^t\mid\theta^t]=\nabla F(\theta^t).\label{eq:federated-inclusion-gradient}
\end{equation}
式\eqref{eq:federated-inclusion-gradient}中的期望等式来自$\E[I_k^t]=\pi_k$。这种逆概率校正不要求各客户端独立入选，但小$\pi_k$会放大方差。常用的参与集合内重新归一化权重是随机比值，通常不能声称对原目标严格无偏。如果网络故障与设备状态相关，又没有可靠参与概率，仅靠增加训练轮数不能恢复缺失工厂的代表性。

\section{联邦优化、隐私与系统约束}
\paragraph{非独立同分布与通信。}
当$P_k(X,Y)$不同，本地更新方向可能冲突，客户端数量、参与率和本地步数会影响收敛。通信压缩、个性化联邦学习和聚类联邦学习分别处理带宽、客户端差异和多群体结构。服务器不能直接看到原始数据，但更新仍可能泄露信息，因此需要安全聚合、差分隐私或加密协议。这里的协议指参与方之间明确规定的消息交换和计算步骤。

\paragraph{评价与边界。}
联邦系统应报告全局性能、最差客户端性能、通信轮数、隐私预算与掉线鲁棒性。联邦学习解决的是数据不能集中这一约束，并不会自动解决标签偏差、恶意客户端或模型公平性问题。

\begin{table}[htbp]
\centering
\caption{联邦学习方法的约束匹配}
\label{tab:federated-method-constraints}
\scriptsize
\setlength{\tabcolsep}{3pt}
\resizebox{\linewidth}{!}{%
\begin{tabular}{p{2.5cm}p{3.8cm}p{3.6cm}p{4.0cm}}
\hline
方法 & 主要修正对象 & 额外代价 & 适用边界 \\
\hline
FedAvg & 基本全局聚合 & 本地漂移未被显式校正 & 分布接近、通信受限的基线 \\
FedProx & 客户端更新偏离 & 近端系数和局部适应之间需调节 & 非 Independent and Identically Distributed (IID) 与算力异构 \\
SCAFFOLD & 系统性梯度漂移 & 控制变量存储和同步 & 客户端长期参与、通信可承受 \\
个性化联邦 & 客户端条件差异 & 共享层与本地层的划分 & 工厂、设备或人群差异明显 \\
安全聚合与 Differential Privacy (DP) & 更新可见性和样本影响 & 密码学开销或噪声导致效用下降 & 需要隐私保证的跨组织协作 \\
\hline
\end{tabular}}
\end{table}

表\ref{tab:federated-method-constraints}将联邦方法与数据异质性、通信及个性化约束对应；比较方法时须先确定主要瓶颈，不能仅按聚合后平均精度选型。
\paragraph{FedAvg的本地更新。}
客户端$k$从全局参数$\theta^t$开始执行$E$步本地更新：
\begin{equation}
 \theta_{k,e+1}^{t}=\theta_{k,e}^{t}-\eta\nabla F_k(\theta_{k,e}^{t}),
 \qquad \theta_{k,0}^{t}=\theta^t.
\end{equation}
服务器聚合$\theta^{t+1}=\sum_kp_k\theta_{k,E}^t$。当各客户端数据同分布且目标函数平滑时，增加本地步数可减少通信；非 IID 情况下，本地漂移会使聚合方向偏离全局梯度。

\paragraph{从本地多步展开看漂移来源。}
先考虑全参与、确定性梯度与共同步长，以免将抽样噪声和漂移混在一起。把$E$步递推求和可得
\begin{align}
 \theta^{t+1}-\theta^t
 &=-\eta\sum_kp_k\sum_{e=0}^{E-1}\nabla F_k(\theta_{k,e}^t)\\
 &=-\eta E\nabla F(\theta^t)-\eta R_t,\\
 R_t&=\sum_kp_k\sum_{e=0}^{E-1}
 [\nabla F_k(\theta_{k,e}^t)-\nabla F_k(\theta^t)].
\end{align}
这里$R_t\in\R^d$就是相对于固定起点梯度的累计偏离。假设各$F_k$梯度为$L$-Lipschitz，且本轮轨迹上$\|\nabla F_k\|\leq G$，由三角不等式有
\begin{equation}
 \|\theta_{k,e}^t-\theta^t\|\leq\eta eG,
 \qquad \|R_t\|\leq L\eta G\frac{E(E-1)}2.
\end{equation}
于是参数误差上界为$L\eta^2GE(E-1)/2$。这不是完整收敛定理，也未区分同分布曲率与异质性；它明确说明：少通信带来的多步本地训练会累积偏离，而不是无代价地模拟集中式梯度。

异质性可用$\zeta^2(\theta)=\sum_kp_k\|\nabla F_k(\theta)-\nabla F(\theta)\|^2$描述。全局最优点上$\nabla F=0$并不要求各$\nabla F_k=0$，因此本地继续优化仍可能把参数拉向不同方向。一个可直接核算的例子是$F_k(\theta)=a_k(\theta-b_k)^2/2$，其中$\theta\in\R$、$a_k>0$。递推给出
\begin{equation}
 \theta_{k,E}=b_k+(1-\eta a_k)^E(\theta^t-b_k).
\end{equation}
当$|1-\eta a_k|<1$且$E\to\infty$，聚合趋于$\sum_kp_kb_k$，但全局最优点是$\sum_kp_ka_kb_k/\sum_kp_ka_k$。除特殊情形外二者不同；本地模型各自训练得更充分，不代表全局模型更接近目标最优点。

\paragraph{隐私的三个层次。}
数据不上传不等于隐私完好。更新可能通过梯度反演暴露样本特征，恶意客户端可能投毒，服务器也可能推断客户端属性。安全聚合保护服务器免于看到单个更新；差分隐私在更新上加入噪声，提供形式化的$(\epsilon,\delta)$保证；鲁棒聚合则针对异常更新，但不能自动提供隐私。

\paragraph{个性化联邦学习。}
全局模型适合共享规律，个性化模型允许客户端保留本地参数：
\begin{equation}
 \theta_k=\theta_{shared}+\theta_{k,local}.
\end{equation}
共享层、个性化头、元学习初始化和客户端聚类是常见路线。工业设备差异很大时，追求一个平均最优模型可能不如学习一组有结构的个性化模型。

\paragraph{个性化目标为什么需要耦合。}
仅写$\theta_k=\theta_{shared}+\theta_{k,local}$还没有定义算法，因为给共享参数加一个向量、给每个本地参数减去同一向量不会改变预测。一个明确的软共享目标是
\begin{equation}
 \min_{\theta,\{v_k\}}\sum_kp_k
 \left[F_k(v_k)+\frac\lambda2\|v_k-\theta\|^2\right],
 \quad \theta,v_k\in\R^d,\quad\lambda>0.
\end{equation}
固定$\theta$，各客户端独立优化其正则化风险；固定$v_k$，对$\theta$求导得到$\lambda\sum_kp_k(\theta-v_k)=0$，故$\theta=\sum_kp_kv_k$。这一交替过程把共享信息理解为各个性化参数的参照，而不是强迫它们完全相同。$\lambda$很大时接近共享模型，很小时接近独立训练，但数据少的客户端可能因此过拟合。这里的$v_k$是长期保留的个性化模型，与 FedProx 中每轮重置的局部变量不同。

若不同工厂的输出标签或传感器结构并不相同，可令共享编码器为$f_\theta$、本地头为$h_{\psi_k}$，优化$\sum_kp_kF_k(\theta,\psi_k)$。此时只能聚合语义和维度一致的共享参数，不能把不同类别顺序的分类头逐坐标平均。对从未参与的新工厂，还必须规定支持样本和本地适应步骤；这个问题将在元学习章进一步展开。

\paragraph{通信和系统约束。}
量化、稀疏化、低秩更新和客户端选择降低通信量。实际系统还面对带宽波动、设备掉线、异构算力和更新延迟。实验应记录每轮参与客户端数、有效样本量、字节数、端侧训练时间和掉线率，而不能只报告最终准确率。

\paragraph{压缩误差与误差反馈。}
设待上传差分$u_k\in\R^d$，压缩器$Q$可返回量化值或稀疏向量。若$\E[Q(u)\mid u]=u$且$\E\|Q(u)-u\|^2\leq\omega\|u\|^2$，压缩不改变条件均值，却增加更新方差；确定性的最大幅度稀疏化通常不满足无偏性。误差反馈保存残差$r_k^t$，令
\begin{equation}
 s_k^t=Q(u_k^t+r_k^t),\qquad
 r_k^{t+1}=u_k^t+r_k^t-s_k^t.
\end{equation}
由此$s_k^t=u_k^t+r_k^t-r_k^{t+1}$，对轮次求和后中间残差相消：$\sum_{t=0}^{T-1}s_k^t=\sum_{t=0}^{T-1}u_k^t+r_k^0-r_k^T$。它解释了为何未发送信息可以延后补偿，但残差有界及收敛仍需压缩器和目标的额外条件。长期掉线、状态丢失或模型版本变化都可能破坏这种补偿；稀疏上传还要计入索引字节，不能只数非零数值。

\section{联邦诊断案例与验收}
将不同工厂作为客户端，训练共享的故障表征和个性化的告警头。一轮 FedAvg 与集中式梯度下降的差异来自本地多步更新和客户端数据权重；安全聚合只能隐藏单个更新，不能防止模型投毒，因此验收还应约束最差客户端性能、异常更新比例和通信失败后的恢复行为。

\paragraph{一轮联邦训练的完整过程。}
服务器首先选择客户端子集并发送当前模型。客户端在本地数据上训练若干步，计算更新后上传参数或差分；服务器按样本量或其他权重聚合，再将新模型发送给下一轮。这个过程的关键不是“把 SGD 分散执行”这么简单，而是每轮参与者、数据量、本地步数和网络状态都会改变优化轨迹。

\begin{figure}[htbp]
\centering
\begin{tikzpicture}[>=Latex, node distance=0.9cm, every node/.style={font=\small}]
 \node[draw, rounded corners, fill=yellow!20, minimum width=2.5cm] (server) {\shortstack{服务器\\全局模型}};
 \node[draw, rounded corners, fill=blue!10, below left=1.2cm and 1.8cm of server, minimum width=2.2cm] (c1) {\shortstack{客户端 1\\本地数据}};
 \node[draw, rounded corners, fill=green!12, below=1.2cm of server, minimum width=2.2cm] (c2) {\shortstack{客户端 2\\本地数据}};
 \node[draw, rounded corners, fill=orange!12, below right=1.2cm and 1.8cm of server, minimum width=2.2cm] (c3) {\shortstack{客户端 3\\本地数据}};
 \draw[<->, thick] (server)--(c1); \draw[<->, thick] (server)--(c2); \draw[<->, thick] (server)--(c3);
 \node[below=0.3cm of c2,font=\small] {双向通信：下发模型 / 上传更新 / 聚合};
\end{tikzpicture}
\caption{联邦学习的服务器--客户端循环；原始数据留在客户端，但更新仍需保护。}
\label{fig:federated-training-round}
\end{figure}
图\ref{fig:federated-training-round}展开一轮下发、本地训练和回传聚合的关系；客户端计算与通信构成不同成本，且回传更新本身仍可能包含敏感信息。

\paragraph{客户端漂移与 FedProx。}
非 IID 数据下，本地目标的最优方向可能不同。客户端执行过多本地步数后，更新会远离全局模型，聚合后产生客户端漂移。FedProx 在本地目标中加入近端项：
\begin{equation}
 \min_\theta F_k(\theta)+\frac{\mu}{2}\|\theta-\theta^t\|_2^2.
\end{equation}
近端系数限制本地模型偏离服务器模型的幅度，适合数据和算力差异较大的系统，但过大的$\mu$会降低本地适应能力。

对该目标求梯度，本地第$e$步为
\begin{equation}
 \theta_{k,e+1}=\theta_{k,e}-\eta
 [\nabla F_k(\theta_{k,e})+\mu(\theta_{k,e}-\theta^t)].
\end{equation}
起点$\theta_{k,0}=\theta^t$处近端梯度为零，所以第一步与 FedAvg 相同；它主要约束之后的偏离。对二次目标$F_k(v)=\tfrac12(v-b_k)^{\mathsf T}A_k(v-b_k)$，其中$A_k\in\R^{d\times d}$对称半正定、$\mu>0$，驻点方程为$(A_k+\mu I)v=A_kb_k+\mu\theta^t$，从而
\begin{equation}
 v-\theta^t=-(A_k+\mu I)^{-1}\nabla F_k(\theta^t).
\end{equation}
沿$A_k$的特征方向，位移被$1/(a+\mu)$缩放，展示了近端项如何抑制偏离。对非凸神经网络，近端项不自动使目标凸；若 Hessian 特征值下界为$-\rho$，只有$\mu>\rho$才足以使该区域的正则化 Hessian 正定。实际内层通常仅近似求解，应同时报告本地停止准则，而非默认获得精确极小点。

\paragraph{控制变量与 SCAFFOLD。}
SCAFFOLD 为服务器和客户端维护控制变量，用于估计并校正本地更新中的系统性偏差。其思想是把“客户端自身的梯度偏移”和“当前模型真正的全局下降方向”区分开。该类方法通常以额外状态和通信为代价改善非 IID 收敛，实验时必须报告控制变量的存储和同步开销。

具体地，设$c,c_k\in\R^d$分别估计全局和本地梯度方向，本地更新改为
\begin{equation}
 v_{k,e+1}=v_{k,e}-\eta(g_k(v_{k,e})-c_k+c),
 \qquad v_{k,0}=\theta^t.
\end{equation}
若在同一参考点$c_k=\nabla F_k(\theta^t)$、$c=\sum_kp_kc_k$，则第一步修正方向的条件期望恰为$\nabla F(\theta^t)$，不再随客户端改变。后续各本地点不同，校正并不意味着始终拥有精确全局梯度。一种固定$E,\eta$的控制变量更新为
\begin{equation}
 c_k^{new}=c_k-c+\frac{\theta^t-v_{k,E}}{E\eta}
 =\frac1E\sum_{e=0}^{E-1}g_k(v_{k,e}).
\end{equation}
第二个等号由本地递推求和得到，说明新控制变量代表本轮沿轨迹的平均梯度。若维护加权关系$c=\sum_kp_kc_k$，服务器应更新$c^{new}=c+\sum_{k\in S_t}p_k(c_k^{new}-c_k)$，未参与者保留旧状态。经典均匀客户端设定取$p_k=1/K$；不同权重、服务器步长和参与方案需要配套实现，不能混用聚合规则。陈旧控制变量以及额外的$d$维上传，是该方法的实际代价。

\paragraph{差分隐私与隐私预算。}
必须先指定保护单位。样本级邻接表示训练集仅相差一条记录；客户端级邻接表示相差一个客户端的全部贡献。单样本梯度裁剪通常用于样本级保护，不能直接称为客户端级保护。随机机制$M$满足$(\epsilon,\delta)$差分隐私，是指任意相邻数据集$D,D'$和可测输出事件$A$都有
\begin{equation}
 \Pr[M(D)\in A]\leq e^\epsilon\Pr[M(D')\in A]+\delta.
\end{equation}
它约束观察结果因一单位数据变化而改变的程度，不是“攻击一定失败”的概率声明。

以客户端级发布为例，先把每个上传向量$u_k\in\R^d$裁剪为
\begin{equation}
 \bar u_k=u_k\min\{1,C/\|u_k\|_2\},\qquad
 M=\frac1m\left(\sum_k\bar u_k+\xi\right),\quad
 \xi\sim\mathcal N(0,\sigma^2C^2I_d),
\end{equation}
其中零向量保持为零，$m$是固定公开归一化常数，缺席贡献以零计。加入或删除一个客户端时，和的$\ell_2$敏感度至多为$C$；替换一个客户端时由三角不等式得到$2C$。因此两种邻接定义不可使用相同敏感度而不作说明。对一次发布、加入删除邻接、$0<\epsilon<1$，经典高斯机制的一个充分条件是$\sigma\geq\sqrt{2\log(1.25/\delta)}/\epsilon$；该式是保守校准条件，不是多轮训练的最终预算。

除以$m$后，敏感度与噪声标准差同时缩小，不能据此声称隐私自动增强。裁剪把过大更新的影响截断，也引入偏差；噪声期望为零，但其平均平方范数为$d\sigma^2C^2/m^2$，高维、少参与者时可能明显损害效用。若服务器先看到未扰动单个更新，再在总和上加噪声，则保证针对发布结果的观察者，而不保护客户端免受该服务器观察；必须明确可信聚合者，或使用安全聚合与按约定步骤进行的分布式噪声生成。

对于自适应的$T$轮发布，基本组合给出$\epsilon_{tot}\leq\sum_t\epsilon_t$、$\delta_{tot}\leq\sum_t\delta_t$。实际可用更紧的隐私会计，但需声明抽样方式、采样率、噪声乘子、邻接定义和会计方法。抽样放大结论依赖抽样独立性及可见信息，不能对按故障分数挑选客户端的过程直接套用均匀抽样公式。

\paragraph{安全聚合与攻击面。}
安全聚合使服务器只能看到聚合结果，不能直接读取单个客户端更新，但它不能防止恶意客户端上传投毒或后门更新。鲁棒聚合可以按坐标中位数、截断均值或相似性筛除异常，却可能误伤来自真实少数群体的正常更新。安全性需要从服务器、客户端、通信通道和模型输出四个方面分别分析。

\paragraph{安全聚合为何能只显露总和。}
下面只用代数说明掩码抵消，不把它当作可直接部署的完整密码协议。设更新经约定量化后为有限域$\mathbb F_q^d$上的向量$u_k$，每对客户端共享随机掩码$r_{kj}$，并约定$r_{jk}=-r_{kj}$。客户端发送
\begin{equation}
 \widetilde u_k=u_k+\sum_{j\ne k}r_{kj},\qquad
 \sum_k\widetilde u_k=\sum_ku_k+
 \sum_{k<j}(r_{kj}+r_{jk})=\sum_ku_k.
\end{equation}
掩码在总和中抵消，单个带掩码消息却不直接等于原更新。实际协议还需要认证、密钥协商、门限恢复及串谋假设；若某客户端掉线，它对应的掩码未必成对抵消，需要恢复机制，且不能同时恢复足以解开其原始贡献的全部秘密。量化范围和模数还要防止解码溢出。

服务器得到总和仍可能通过参与集合差分或很小的聚合组推断贡献，因此还需最小聚合规模、查询限制及必要的差分隐私。另一个根本冲突是：坐标中位数不是总和的函数，服务器若只能获得总和，就不能直接执行基于单个更新的筛选。安全聚合与鲁棒聚合必须联合设计，而不是把两个算法名称并列便认为两项保证同时成立。

\paragraph{掉线、异构与公平性。}
慢客户端会拖延同步，掉线客户端可能造成样本代表性变化；按样本量加权又可能让大客户端压过小群体。评价时应报告平均性能、最差客户端、不同客户端规模区间和掉线条件下的性能。对于工业多工厂系统，公平性不仅是统计指标，也关系到不同工厂是否都能获得足够的故障检测能力。

在多工厂轴承诊断中，可把工厂、产线或设备群定义为客户端。首先固定一个不参与训练的跨工厂测试集，并在每轮记录参与客户端、样本量、本地步数、上传字节数和训练时长。随后构造三种非 IID 情形：不同工厂故障比例不同、不同工厂只观测部分类别、不同工厂的相机或传感器分布不同。FedAvg、FedProx、SCAFFOLD 和个性化头应使用相同总通信预算比较，同时报告全局平均、宏平均、最差工厂和各类故障召回。

隐私验收不能只检查原始文件是否离开工厂。应分别测试更新反演、成员推断、恶意客户端投毒和服务器串谋；对差分隐私记录每轮裁剪、噪声和累计$(\epsilon,\delta)$，对安全聚合记录阈值、掉线恢复和密钥失效行为。若加入噪声后小工厂的少数故障召回显著下降，平均准确率可能仍然稳定，但系统已经产生了不公平的风险分配，因此最差客户端和少数类指标必须进入停止条件。

联邦优化、客户端漂移和隐私边界可结合下面两份非论文材料学习：一份是框架实践，另一份是隐私机制说明。
\section*{辅助学习材料}
\begingroup
\emergencystretch=3em
\begin{itemize}
\item \textbf{Quickstart PyTorch}；作者/平台：Flower；英文；官方框架教程。阅读 ClientApp、ServerApp、FedAvg 与本地模拟流程，观察数据分区、轮次聚合和客户端评估。\href{https://flower.ai/docs/framework/tutorial-quickstart-pytorch.html}{在线阅读}。
\item \textbf{TensorFlow Privacy：隐私保护指南}；作者/平台：TensorFlow；中文；官方指南。阅读 DP-SGD 的介绍，理解限制单个样本影响的动机及其与联邦学习用户级隐私的区别；隐私预算还需结合本章的会计分析。\href{https://tensorflow.google.cn/responsible\_ai/privacy}{在线阅读}。
\item \textbf{详解联邦学习Federated Learning}；知乎；中文解说。画出本地训练与服务器聚合流程，检查客户端权重；标题与主题据必应索引，原页受限。\href{https://zhuanlan.zhihu.com/p/79284686}{在线阅读}。
\item \textbf{终于有人把联邦学习讲明白了}；CSDN；中文博客。对照非独立同分布和通信约束；数据不出本地不等于已经满足隐私保证；标题与主题据必应索引，原页受限。\href{https://blog.csdn.net/hzbooks/article/details/118097998}{在线阅读}。
\item \textbf{《联邦学习技术介绍、应用和FATE开源框架》第1课！了解一下AI的新风口？}；上传者：FedAI联邦学习；bilibili；中文课程。对照本章联邦协作流程与应用边界，区分框架实现和隐私保证；旧课程的软件接口以当前官方文档为准。2026年10月4日公开元数据播放12,741次，高于本次异步联邦教程候选的6,701次；已核对公开课程简介、标题和计数，未逐帧观看。\href{https://www.bilibili.com/video/BV1JE411p72s/}{观看视频}。
\end{itemize}
\endgroup
\paragraph{本章小结。}
联邦学习要求数据留在各工厂，但这不会替我们决定应该最小化哪一种风险。所有工厂各做一步、再按样本聚合时，FedAvg 可还原集中式梯度；每个工厂本地走多步后，各自停在不同位置，聚合方向就会产生偏差。FedProx 用近端项限制本地模型离开全局模型太远，SCAFFOLD 用控制变量修正客户端漂移，个性化方法则允许不同工厂保留不同的预测头。它们解决的不是同一个优化问题，不能只按通信轮数或平均准确率排顺序。

隐私与可靠性同样必须拆开：安全聚合限制单个更新的可见性，差分隐私限制单个保护单位对发布分布的影响，鲁棒聚合处理异常贡献。有效系统还要在掉线、少数故障和最差工厂上通过验收。本章的基本立场不是尽可能多地分散计算，而是在可说明的统计目标、信任边界与资源预算之间选择可执行的协作方式。

能够共同训练以后，还要检查怎样衡量样本相似性。高维读数距离近，可能只是温度量纲较大，不一定代表故障相似。下一章的流形学习研究局部邻近关系能否反映转速、负载等少量变化因素。理解这些关系，也有助于判断联邦学习中所谓相似工厂究竟是工况相似，还是仅仅表示中的数值接近。
